% Intake and Shooter Control Changes -- Drive Team Brief
% Build: latexmk -pdf intake-control-changes.tex   (then latexmk -c to remove build files)

\documentclass[11pt]{article}

\usepackage[T1]{fontenc}
\usepackage[utf8]{inputenc}
\usepackage{lmodern}
\usepackage[margin=0.9in]{geometry}
\usepackage{booktabs}
\usepackage{tabularx}
\usepackage{xltabular}
\usepackage{array}
\usepackage{enumitem}
\usepackage{amssymb}
\usepackage[table]{xcolor}
\usepackage{titlesec}
\usepackage[most]{tcolorbox}
\usepackage[colorlinks=true,linkcolor=black,urlcolor=blue]{hyperref}

\definecolor{accent}{HTML}{1F4E79}
\definecolor{notebg}{HTML}{FFF4D6}
\definecolor{noteframe}{HTML}{C98A00}
\newtcolorbox{drivenote}[1]{enhanced, breakable, colback=notebg, colframe=noteframe,
  boxrule=0.8pt, arc=2pt, left=6pt, right=6pt, top=4pt, bottom=4pt,
  fonttitle=\bfseries, title={#1}}
\titleformat{\section}{\large\bfseries\color{accent}}{}{0pt}{}
\titlespacing*{\section}{0pt}{14pt}{4pt}

\setlength{\parindent}{0pt}
\setlength{\parskip}{5pt}
\setlength{\LTpre}{4pt}
\setlength{\LTpost}{6pt}
\renewcommand{\arraystretch}{1.25}
\newcolumntype{L}{>{\raggedright\arraybackslash}X}
\newcommand{\dash}[1]{\texttt{#1}}
\setlist{itemsep=2pt, topsep=3pt}
\newlist{checklist}{itemize}{1}
\setlist[checklist]{label=$\square$, itemsep=2pt, topsep=3pt}

\begin{document}

{\LARGE\bfseries\color{accent} Intake and Shooter Control Changes --- Drive Team Brief}\\[4pt]
{\small 23 November 2026 \quad$\cdot$\quad Branch \texttt{quadecnm/intake\_sensors\_test}}

\section*{Summary}

One press now moves the intake all the way out or all the way in, and two new limit sensors
stop it. The code checks the two lead screws against each other every 10\,ms and stops both if
either jams. The intake roller won't spin until the intake is 75\% out, so its wheels can't rub
on internal parts.

The right trigger now aims at the middle of our hub and picks the shooter speed from the
distance, but only inside our ALLIANCE ZONE, where scoring is legal (rule G407). If the Limelight
can't see a hub tag, the driver aims and the speed comes from the robot's tracked field position.
Outside our zone, or when the robot can't tell where it is, it is the old fixed 3000 RPM shot
with the driver aiming (a pass). In every mode it feeds by itself once the wheels are at speed.

It has only run in simulation so far; see the last section for what still needs testing on the
robot.

\section*{Controls: old vs new}

The extension went from hold-to-move to press-once, and retract moved from D-pad right to D-pad
up.

\begin{xltabular}{\linewidth}{>{\raggedright\arraybackslash}p{2.6cm} L L}
\toprule
\textbf{Control} & \textbf{Old (main branch)} & \textbf{New} \\
\midrule
\endhead
D-pad down    & Hold to extend; stops when released  & Press once: extends until fully out (S2 sensor) \\
D-pad right   & Hold to retract; stops when released & Not used \\
D-pad up      & Not used                             & Press once: retracts until fully in (S1 sensor) \\
Left trigger  & Intake roller in                     & Intake roller in, but only spins once the intake is 75\% out \\
A             & Reverse intake roller                & Same, also waits for 75\% out \\
B             & Storage roller + indexer             & No change \\
X             & Reverse indexer                      & No change \\
Right trigger & Shooter at a fixed 3000 RPM          & In our zone: aims at the hub's middle with RPM from distance (or driver aims, RPM from tracked position). Outside our zone: fixed 3000 RPM, driver aims. Feeds once at speed; intake stows during the shot \\
\bottomrule
\end{xltabular}

\begin{itemize}
  \item \textbf{Which controller:} mechanism buttons go on the operator controller if both
    controllers are connected when the robot boots; otherwise on the driver controller.
    Unchanged from before.
  \item \textbf{Right trigger:} see the next section.
\end{itemize}

\section*{Right trigger: auto-aim, position, or fixed}

Hold the right trigger to shoot. Every 20\,ms the robot picks one of three modes on its own:

\begin{xltabular}{\linewidth}{>{\raggedright\arraybackslash}p{2.3cm} L L L}
\toprule
 & \textbf{AUTO\_AIM} & \textbf{POSE} & \textbf{FIXED} \\
\midrule
\endhead
When          & In our alliance zone, a hub tag seen in the last quarter second & In our alliance zone, no hub tag, robot's position trusted & Outside our zone (passing), or the robot can't tell where it is \\
Turning       & Robot turns itself to the hub's middle; right stick ignored & Right stick; driver aims & Right stick; driver aims \\
Shooter speed & From the measured distance to the hub's middle & From the tracked position to the hub's middle & Fixed 3000 RPM \\
Driving       & Left stick; left bumper = slow mode & Same & Same \\
Feeding       & Automatic, once within 100 RPM of the target & Same & Same \\
Intake        & Stows at the start of the shot & Same & Same \\
\bottomrule
\end{xltabular}

\begin{itemize}
  \item \textbf{Aim point:} the middle of the hub. All eight tags on our hub count, two per face.
    Each tag in view gives an estimate of where the hub's middle is, and they are averaged, so the
    aim does not jump between tags. The camera's off-center mount is accounted for. The other
    alliance's hub is never a target.
  \item \textbf{Close range:} with the Limelight mounted level (as currently set in code), it only
    sees a hub tag from about 3\,m or more from the hub's middle. Closer, POSE mode still picks the
    right speed, but the driver aims. Tilting the Limelight up about 25$^\circ$ (a reprinted mount)
    would let auto-aim work up close.
  \item \textbf{Why POSE exists:} without it, a close shot with no tag in view would run the fixed
    3000 RPM and overshoot. When the robot's position is trusted is in the next section.
  \item \textbf{Switching:} as soon as a hub tag is seen (in our zone), aiming takes over. Feeding
    pauses until the wheels reach whichever speed is current.
  \item \textbf{Old method:} on the main branch the right trigger only spun the shooter at 3000 RPM
    and the driver fed with B. Now feeding is automatic, and B does not feed while the right
    trigger is held.
  \item \textbf{Release:} letting go stops the shooter and feeder. What the intake does next is
    below.
  \item \textbf{Two controllers:} the operator holds the trigger, but in POSE and FIXED the
    driver's right stick aims. Call out ``no tag'' or ``passing'' so the driver knows to aim.
  \item \textbf{Dashboard:} \dash{Score/Mode} shows AUTO\_AIM, POSE or FIXED;
    \dash{Score/InAllianceZone} shows whether the robot thinks it may score.
\end{itemize}

\begin{drivenote}{Alliance zone: scoring only from our side (rule G407)}
The game manual (section 5.3) makes each ALLIANCE ZONE 158.6\,in (about 4.03\,m) deep from that
alliance's wall, across the full field width; the hub's alliance-wall face sits on its edge.
Rule G407 (MAJOR FOUL): launching into our hub is only legal with our bumpers partly or fully in
our alliance zone.
\begin{itemize}[leftmargin=*]
  \item The robot counts as in its zone when its \emph{center} is inside, leaving about half a robot
    of margin for position error.
  \item With a hub tag in view, the robot works out where it is from the tag itself; otherwise from
    its trusted tracked position. If it can't tell, it uses FIXED.
  \item In the neutral zone or the other alliance's zone, every shot is FIXED at 3000 RPM with the
    driver aiming, even with a hub tag (ours or theirs) clearly in view. That is the passing shot.
\end{itemize}
\end{drivenote}

\medskip
\begin{drivenote}{Tuning the shot: measure to the middle of the hub}
Shot distances are now measured from the robot's center to the \textbf{middle of the hub}, not to
the tag on its face. The two placeholder tune points were moved out by the face-to-middle distance
(about 0.6\,m) so the same spots on the floor still get the same speed. When tuning, measure each
shot's distance to the hub's middle (or read \dash{Score/PoseDistanceMeters} /
\dash{Vision/HubDistanceMeters}) and record the RPM that scored.
\end{drivenote}

\medskip
\begin{drivenote}{Intake during a right-trigger shot}
Pulling the trigger always starts stowing the intake. Whether it goes back out on release depends
only on where it was \emph{when the trigger was pulled}:

\smallskip
\begin{tabularx}{\linewidth}{>{\raggedright\arraybackslash}p{5.2cm} L}
\toprule
\textbf{Intake when the trigger is pulled} & \textbf{After the trigger is released} \\
\midrule
Fully out (S2 closed)            & \textbf{Goes back out to S2} (about 4.3\,s) \\
Still extending toward S2        & \textbf{Goes back out to S2} \\
Partly out, stopped              & Stays wherever the stow got to \\
Retracted, or retracting         & Stays in \\
\bottomrule
\end{tabularx}

\smallskip
\begin{itemize}[leftmargin=*]
  \item \textbf{Feeding does not wait for the stow.} The first ball feeds about 0.4\,s after the
    pull, while the intake is still moving in (a full stow takes about 4.2\,s).
  \item \textbf{Quick release:} let go before the stow finishes and the intake reverses straight
    back out (if it started fully out).
  \item \textbf{Intake buttons are ignored while the trigger is held.} To keep the intake in after
    a shot that started fully out, press D-pad up after releasing.
  \item \textbf{Intake roller:} if left trigger is still held, the roller restarts once the intake
    passes 75\% on the way back out.
  \item \textbf{No re-extend} if the intake is not homed, or if the robot is disabled mid-shot.
\end{itemize}
\end{drivenote}

\medskip
Simulation of the right trigger in teleop (camera level, as currently set in code):

\begin{xltabular}{\linewidth}{>{\raggedright\arraybackslash}p{5.0cm} L >{\raggedright\arraybackslash}p{2.6cm} L}
\toprule
\textbf{Where the robot is} & \textbf{Camera saw} & \textbf{Mode} & \textbf{Result} \\
\midrule
\endhead
Our zone, facing away from the hub & 2 wall tags & POSE, 3009 RPM & right stick turned it; fed at 0.41\,s \\
Mid-field, no tags for 11\,s & nothing & FIXED, 3000 RPM & right stick turned it; fed at 0.42\,s \\
Our zone, facing the hub 10$^\circ$ off & 10 tags, ours 25 and 26 & AUTO\_AIM, 3003 RPM & turned itself to 2.4$^\circ$ off the hub's middle; fed at 0.41\,s \\
Blue robot in red's zone, facing red's hub & 10 tags, red's hub & FIXED, 3000 RPM & right stick turned it; fed at 0.41\,s \\
Blue robot in red's zone, passing back & 10 tags, incl.\ our hub's tag 18 from 9.5\,m & FIXED, 3000 RPM & right stick turned it; fed at 0.41\,s \\
Red robot in blue's zone, facing blue's hub & 10 tags, blue's hub & FIXED, 3000 RPM & right stick turned it; fed at 0.43\,s \\
\bottomrule
\end{xltabular}

In every case the intake started in and stayed in after release. With the intake fully out before
the shot, it stowed to S1 in 4.18\,s while the trigger was held and went back out to S2 in 4.18\,s
after release.

Auto-aim from around the hub in simulation (distances to the hub's middle):

\begin{xltabular}{\linewidth}{L >{\raggedright\arraybackslash}p{2.0cm} r r >{\raggedright\arraybackslash}p{2.3cm} r}
\toprule
\textbf{Robot position} & \textbf{Hub tags used} & \textbf{True m} & \textbf{Meas. m} & \textbf{Mode} & \textbf{Aim err.} \\
\midrule
\endhead
Our zone, close tune point           & none       & 2.10 & --   & POSE (2201 RPM) & driver \\
Our zone                             & none       & 2.60 & --   & POSE (2586 RPM) & driver \\
Our zone                             & 25, 26     & 3.10 & 3.08 & AUTO\_AIM & 0.0$^\circ$ \\
Our zone                             & 25, 26     & 3.60 & 3.55 & AUTO\_AIM & 0.0$^\circ$ \\
Our zone, far tune point             & 25, 26     & 4.10 & 4.09 & AUTO\_AIM & 0.0$^\circ$ \\
Our zone, turned 20$^\circ$          & 25, 26     & 3.10 & 3.13 & AUTO\_AIM & 1.8$^\circ$ \\
Our zone, 1\,m to the side           & 21, 24, 25, 26 & 3.26 & 3.25 & AUTO\_AIM & 0.3$^\circ$ \\
Beside the hub / center-field side (outside our zone) & 18--27 & 3.10 & -- & FIXED & driver \\
Red: own zone / beside its hub       & 9, 10 / 2, 11 & 3.10 & 3.09 / -- & AUTO\_AIM / FIXED & 0.0$^\circ$ / driver \\
\bottomrule
\end{xltabular}

Practice-field simulation: only one hub's tags, robot powered up believing it was 6--16\,m and
30--150$^\circ$ from where it really was, then driving in to the close tune point:

\begin{xltabular}{\linewidth}{L >{\raggedright\arraybackslash}p{3.0cm} >{\raggedright\arraybackslash}p{3.2cm}}
\toprule
\textbf{Case} & \textbf{Close shot} & \textbf{Driver's forward} \\
\midrule
\endhead
Blue, looked at the hub while disabled & POSE 2147 RPM (2143 needed) & never changed \\
Blue, enabled right away & POSE 2158 RPM (2149 needed) & at most 1$^\circ$ per 20\,ms; back to normal \\
Red, looked at the hub while disabled & POSE 2152 RPM (2146 needed) & never changed \\
Red, enabled right away (150$^\circ$ off) & POSE 2152 RPM (2146 needed) & at most 1$^\circ$ per 20\,ms; back to normal \\
Blue-tag hub, Driver Station on red & FIXED 3000 RPM (not in red's zone) & never changed \\
\bottomrule
\end{xltabular}

The simulation uses an ideal camera and field; it does not prove the real Limelight is set up
right. See the checklist at the end.

\section*{Robot position: when it is trusted, and practice fields}

POSE mode and the alliance-zone check rely on the robot knowing where it is. It tracks that with
the wheels, and corrects it from any AprilTag the Limelight sees (any tag on the field, any
alliance). At power-up it thinks it is at a corner of the field facing 0$^\circ$; in a match, the
auto routine sets the right position at the start.

\begin{itemize}
  \item \textbf{Trusted:} a camera position agrees with the tracked one within 0.25\,m and
    5$^\circ$. It stays trusted for 10\,s after the last agreeing picture.
  \item \textbf{Not trusted:} after power-up, after pressing Y, after a camera picture disagrees by
    more than 0.5\,m or 15$^\circ$, or 10\,s without an agreeing picture. Close shots then use FIXED
    instead of guessing.
  \item \textbf{Snap:} when the position is not trusted and the Limelight sees two or more tags at
    once that disagree by more than 0.5\,m or 15$^\circ$, the robot jumps straight to the camera's
    position. This also covers a robot code restart in the middle of a match.
  \item \textbf{The driver's ``forward'' never jumps.} If the snap happens while enabled, joystick
    ``forward'' first keeps pointing the same way in the room, then eases back to the normal
    alliance direction at a fixed 36$^\circ$ per second: a few degrees settle before anyone
    notices; a half-turn takes 5\,s and feels like a slow drift (a straight push can bend into a U)
    to steer through. While disabled it goes straight back to normal.
  \item \textbf{Y button} (reset field-centric heading): makes the robot's current facing
    ``forward'', stops any easing, and stops trusting the position until a picture agrees again.
  \item \textbf{Picking the alliance:} the robot only knows its alliance from the Driver Station
    (the field system in a match; the \textbf{Team Station} drop-down in practice). Without one, it
    uses FIXED.
\end{itemize}

\begin{drivenote}{Practicing with only a hub set up}
\begin{itemize}[leftmargin=*]
  \item Tag the practice hub with one alliance's tag IDs from the 2026 field map, on the right
    faces and at the right height and spacing: blue 18--21 and 24--27, red 2--5 and 8--11. Two tags
    per face.
  \item Set the Driver Station Team Station to that same alliance. (A blue-tag hub with the
    Driver Station on red simply gives FIXED shots, because the robot is not in the red zone.)
  \item Before enabling, let the robot sit for a second or two where the Limelight can see a face
    of the hub with both tags. It snaps to the right position while disabled, and close shots use
    POSE from then on.
  \item Enabling without that also works: the robot snaps as soon as it sees two hub tags, and
    ``forward'' eases back as above.
\end{itemize}
\end{drivenote}

\medskip
\section*{How the extension works now}

The intake now runs itself to a sensor after one press, instead of chasing fixed motor
positions while a button is held.

\begin{xltabular}{\linewidth}{>{\raggedright\arraybackslash}p{2.8cm} L L}
\toprule
 & \textbf{Old} & \textbf{New} \\
\midrule
\endhead
Starting a move    & Hold the button the whole time                   & One press; the move finishes on its own \\
Where it stops     & Fixed motor positions, guessed from power-up     & The S1 (in) or S2 (out) limit sensor \\
Two screws         & Right screw just copied the left                 & Each screw checked on its own; the faster one waits for the slower \\
Speed              & Motor's position control, untuned                & 15 motor turns/s = 3\,in/s; full extend is about 4.3\,s \\
Changing direction & Release one button, hold the other               & Press the other direction mid-move and it reverses \\
Disable            & Screws stop                                      & Screws stop and the old move does not resume on enable \\
\bottomrule
\end{xltabular}

The nut moves 0.2\,in per screw turn (1/2-5 ACME; the motors drive the screws directly). Sensor
to sensor is 11.75\,in, about 59 motor turns.

\section*{New sensors}

A CTRE CANdi with two limit sensors replaces the old single limit switch on the roboRIO (DIO 0).

\begin{xltabular}{\linewidth}{>{\raggedright\arraybackslash}p{3.4cm} >{\raggedright\arraybackslash}p{3.6cm} L}
\toprule
\textbf{Sensor} & \textbf{Where} & \textbf{What the code does with it} \\
\midrule
\endhead
CANdi S1 & Intake fully retracted & Ends a retract; zeroes both screw encoders; marks the intake ``homed'' \\
CANdi S2 & Intake fully extended, 11.75\,in from S1 & Ends an extend; always allows the intake roller \\
Kraken X60 encoders (both screws) & Built into the screw motors & Tracks position between sensors and checks the two screws stay together \\
Limelight (PhotonVision) & Camera on the robot & Finds any of our hub's eight AprilTags to aim at the hub's middle and measure distance; any tag on the field also corrects the robot's tracked position (used by POSE mode) \\
\bottomrule
\end{xltabular}

The CANdi is CAN ID 11 on the roboRIO CAN bus. The hard stops sit about 1/8\,in past each
sensor.

\section*{New safety precautions}

None of these existed before; the old code had no jam detection, no current limits and no
roller interlock.

\begin{xltabular}{\linewidth}{>{\raggedright\arraybackslash}p{3.0cm} L L}
\toprule
\textbf{Precaution} & \textbf{Trips when} & \textbf{What happens} \\
\midrule
\endhead
Screws out of step & The two screws are more than 1 turn (0.2\,in) apart & Both stop (fault). The next extend or retract press lines them up first, then moves \\
Screw off target & Either screw is more than 1.5 turns behind where it was told to be & Both stop (fault), same recovery \\
Stuck & Neither screw makes progress for 0.5\,s & Both stop (fault), same recovery \\
Current limits & Always on & Each screw is capped at 40\,A (torque) and 30\,A (battery), so a jam can't wind the actuator up \\
Missed sensor & A move goes 0.1\,in past where its sensor should have triggered & Stops short of the hard stop; Driver Station warning \\
Move timeout & One move runs longer than 12\,s & Stops; Driver Station error \\
Not homed & Powered up away from S1 & Extend blocked; retract moves at most 1.4\,in per press (see next section) \\
Intake roller interlock & Intake less than 75\% out (and S2 not closed) & Intake roller held off; starts by itself at 75\% \\
Lost sensor data & Encoder or CANdi readings fail & Screws stop and the intake roller is held off \\
Disable / start of teleop & Robot disabled, or teleop starts & Extension stops; shooter, indexer and rollers stop; an unfinished auto shot is forgotten \\
\bottomrule
\end{xltabular}

A fault stays until someone presses extend or retract; disabling and re-enabling does not clear
it.

\section*{Power-up and pit procedure}

Power the robot off with the intake fully retracted whenever you can; then it boots ready and
there is nothing to do.

The robot only knows where the intake is after it has seen the retracted sensor (S1). The lead
screws can't be pushed by hand, so if it powers up partly out, walk it in with the motors:

\begin{enumerate}
  \item Check the dashboard: \dash{Intake/Homed} is false and the Driver Station shows ``Intake
    not homed''. Extend (D-pad down) is ignored in this state.
  \item Enable the robot in the pit (tethered).
  \item Press D-pad up. Each press moves the intake in at most 1.4\,in, then stops. Keep
    pressing; from fully out it takes about 9 presses.
  \item When S1 triggers, the intake stops, \dash{Intake/Homed} turns true and the positions read
    0. Everything is normal from then on.
  \item Before queuing, confirm \dash{Intake/Homed} and \dash{Intake/RetractedLimit} are both
    true.
\end{enumerate}

If S1 is broken, step 3 never moves more than 1.4\,in per press, so a dead sensor can't drive it
far. If the intake looks fully in but \dash{Intake/RetractedLimit} stays false, stop pressing and
check the sensor.

\section*{Auto and shooting changes}

Autos now fully deploy the intake before collecting; before, the intake never got all the way
out in auto.

\begin{itemize}
  \item \textbf{Old auto deploy:} the deploy step actually ran the retract code, and a stop 1\,s
    later cut it off anyway.
  \item \textbf{New auto deploy:} the robot waits after its first path until the intake reaches
    S2, about 4.3\,s, then continues.
  \item \textbf{Shots stow the intake:} a right-trigger shot, and every auto shot, retracts the
    intake first. When the trigger is released, it goes back out only if it was fully out when
    the shot started.
  \item \textbf{Intake roller during a shot:} it stops as the intake retracts past 75\% and
    restarts if the intake goes back out.
  \item \textbf{Auto shots are fixed-speed:} the auto's shooter step spins the shooter at 3000 RPM
    and feeds on the auto's own timing. Auto-aim and POSE are right-trigger (teleop) features.
  \item \textbf{Vision at the start of an auto:} the auto sets the robot's position to the path's
    starting point. If the robot was actually placed more than 0.5\,m or 15$^\circ$ from there and
    the Limelight sees two or more tags, the robot snaps to where it really is, and the paths drive
    from there. Smaller placement errors are nudged out gradually. Place the robot on the start
    position as accurately as possible either way.
\end{itemize}

Auto times in simulation, against the 20\,s limit:

\begin{xltabular}{\linewidth}{L r r}
\toprule
\textbf{Auto} & \textbf{Time} & \textbf{Spare} \\
\midrule
\endhead
Rightstart to Mid collcect then shoot                        & 18.80\,s & 1.2\,s \\
RightStart to Mid collect then pass (USE THIS ONE PLAYOFFS)  & 17.08\,s & 2.9\,s \\
Left start to Mid that collects and pass to teamside         & 14.41\,s & 5.6\,s \\
Middle start shoot                                           &  8.64\,s & 11.4\,s \\
\bottomrule
\end{xltabular}

The real robot may run slower than the simulation, so watch the first auto closely.

\section*{Dashboard and Driver Station messages}

These are the new values to watch on the dashboard.

\begin{xltabular}{\linewidth}{>{\raggedright\arraybackslash}p{6.2cm} L}
\toprule
\textbf{Dashboard value} & \textbf{Means} \\
\midrule
\endhead
\dash{Intake/State} & REST, EXTENDING, RETRACTING, REALIGNING (lining the screws up) or FAULT \\
\dash{Intake/Fault} & Why it stopped: SCREWS\_APART, LEFT\_OFF\_TARGET, RIGHT\_OFF\_TARGET or NO\_PROGRESS \\
\dash{Intake/Homed} & True once S1 has been seen since power-up \\
\dash{Intake/RetractedLimit} / \dash{Intake/ExtendedLimit} & S1 / S2 sensor closed \\
\dash{Intake/RollerAllowed} & Intake roller is allowed to spin (75\% out or S2 closed) \\
\dash{Intake/ExtendedPercent} & How far out the intake is \\
\dash{Intake/LeftPos}, \dash{Intake/RightPos}, \dash{Intake/Apart} & Each screw's position in motor turns, and the gap between them \\
\dash{Intake/SensorsOK} & False means the encoder or CANdi readings failed \\
\dash{Score/TargetRpm}, \dash{Score/ShooterRpm}, \dash{Score/Feeding} & Shot RPM wanted vs measured, and whether a ball is being fed \\
\dash{Score/Mode} & AUTO\_AIM, POSE or FIXED \\
\dash{Score/InAllianceZone} & The robot thinks it is in our alliance zone (may score) \\
\dash{Drive/ForwardOffsetDegrees} & How far joystick ``forward'' still is from normal after a snap while enabled \\
\dash{Vision/PoseResets}, \dash{Vision/PoseOffMeters} & Number of snaps; how far the last camera position was from the tracked one \\
\dash{Score/PoseDistanceMeters} & Distance from the robot's tracked position to the hub's middle \\
\dash{Vision/Connected} & False: Limelight off, unplugged, or camera name mismatch \\
\dash{Vision/HubStatus}, \dash{Vision/HubDistanceMeters} & Whether a hub tag is in view (or why not), and the measured distance to the hub's middle \\
\bottomrule
\end{xltabular}

\begin{xltabular}{\linewidth}{>{\raggedright\arraybackslash}p{6.2cm} L}
\toprule
\textbf{Driver Station message} & \textbf{What to do} \\
\midrule
\endhead
``Intake screws stopped, \ldots'' & A jam or the screws got out of step. Clear the jam, then press extend or retract; it lines the screws up first \\
``Intake not homed \ldots'' & Walk the intake in with D-pad up (Power-up and pit procedure) \\
``Intake not homed: extend ignored \ldots'' & Same; extend works after S1 is reached \\
``Intake reached the travel limit without its CANdi sensor triggering \ldots'' & A sensor missed. Check its wiring before the next match \\
``Intake move timed out \ldots'' & A move took over 12\,s. Check for binding \\
\bottomrule
\end{xltabular}

\section*{Still to test on the robot}

Everything above has run only in simulation. Check these with the robot on blocks before a
match:

\begin{checklist}
  \item S1 and S2 read correctly on the dashboard (closed when the intake is at that end, not
    reversed)
  \item \dash{Intake/LeftPos} and \dash{Intake/RightPos} both go up while extending
  \item About 5 motor turns per inch, and about 59 at S2
  \item Several extends and retracts at full speed with no ``screws stopped'' fault; actuator is
    rated for 3\,in/s
  \item The screws coast less than 1/8\,in past each sensor (the hard stop)
  \item Power up partly out, then walk it in with D-pad up; repeat with S1 unplugged and confirm
    each press moves at most 1.4\,in
  \item Hold left trigger during an extend: the roller starts at about 75\% out and stops when
    the intake retracts past it
  \item Limelight runs PhotonVision (not the stock Limelight software), the camera is named
    \texttt{Camera\_\allowbreak Module\_\allowbreak v1}, and \dash{Vision/Connected} is true
  \item Pipeline 0 is an AprilTag pipeline with 3D mode and multi-tag on, using the 2026 field map
    that matches the event (welded or AndyMark); the camera is calibrated at that resolution
  \item Measure the Limelight mount (forward, left, height, tilt) and update the camera constants;
    decide on a 25$^\circ$ up-tilt mount
  \item Shot RPM: the two distance/RPM pairs are placeholders. Tune them on the field, measuring
    to the hub's middle
  \item Shooter reaches about 3000 RPM and holds it (new speed feedforward; assumes Kraken X60
    shooter motors)
  \item Turn away from the hub near it, hold the right trigger: \dash{Score/Mode} reads POSE and
    the RPM matches the distance; the right stick turns the robot
  \item Drive into the neutral zone and the other alliance's zone and hold the right trigger.
    Expect FIXED at 3000 RPM and no auto-aim, even facing a hub, and
    \dash{Score/InAllianceZone} false
  \item Drive across the starting line: \dash{Score/InAllianceZone} should flip as the robot's
    center crosses it
  \item Power up away from the hub and enable right away: once it sees two hub tags it snaps
    (\dash{Vision/PoseResets} counts up) and ``forward'' eases back without jumping
  \item Unplug the Limelight for over 10\,s, hold the right trigger: FIXED at 3000 RPM
  \item Turn back toward the hub mid-shot: it switches to AUTO\_AIM and turns onto the hub's
    middle, from each side of the hub
  \item Extend fully, hold the right trigger until the intake is in, release: it goes back out to
    S2. Repeat starting retracted: it stays in
  \item Watch a ball feed while the intake is still stowing; confirm nothing catches
  \item Check shots land centered; if they land consistently to one side, adjust the aim offset
    constant
  \item Run each auto on the real robot and time it, especially ``Rightstart to Mid collcect then
    shoot''
\end{checklist}

\end{document}
